\subsection*{Partie A -- Étude d'un cercle}
	\begin{enumerate}
		\item Un point {M~$(x;y)$} appartient à $\mathcal{C}$ si et seulement si la distance AM (donc la norme du vecteur $\vectt{AM}$) est égale à {$\dfrac{\sqrt{37}}{2}$.}

		Il faut donc que {$\left\|\overrightarrow{\text{AM}}\right\|^2=\left(\dfrac{\sqrt{37}}{2}\right)^2=\dfrac{37}{4}$}

		Les coordonnées de {$\vectt{AM}$} sont {$\left(\begin{array}{c}x-3\\ y-1\end{array}\right)$} et donc {$\left\|\overrightarrow{\text{AM}}\right\|^2=(x-3)^2+(y-1)^2$.}

		Une équation cartésienne du cercle {$\mathcal{C}$} est donc {$(x-3)^2+(y-1)^2=\dfrac{37}{4}$.}

		\item On vérifie si les coordonnées de B vérifient cette équation :\quad  {$(6-3)^2+(\np{0.5}-1)^2=3^2+(-\np{0.5})^2=9+\np{0.25}=\np{9.25}$}.

		Or \quad {$\np{9.25}=\dfrac{37}{4}$}, \quad donc B appartient à {$\mathcal{C}$.}
	\end{enumerate}

\subsection*{Partie B -- Étude d'un cas particulier}

	\begin{enumerate}
		\item L'équation de {$\mathcal{H}$} est {$y=\dfrac 3 x$}, donc pour {$x=6$}, on obtient {$y=\dfrac 3 6=\np{0.5}$}, qui est bien l'ordonnée de B, qui appartient donc à {$\mathcal{H}$.}

		\item \begin{enumerate}
			\item Les coordonnées de {$\vectt{AD}$} sont {$\left(\begin{array}{c}1\\ 1\end{array}\right)$}, et celles de {$\vectt{AB}$} sont {$\left(\begin{array}{c}3\\ -\np{0.5}\end{array}\right)$.}

			\item On a donc {$\vectt{AD}\cdot\vectt{AB}=1\times3+1\times(-\np{0.5})=\np{2.5}$.}

			\item L'angle  $\widehat{\text{DAB}}$ n'est pas un angle droit, en effet s'il l'était, le produit scalaire {$\vectt{AD}\cdot\vectt{AB}$} serait nul.
		\end{enumerate}
	\end{enumerate}

\subsection*{Partie C -- Étude du cas général}

	\begin{enumerate}
		\item On calcule \quad {$\Delta=(-4)^2-4\times1\times3=4=2^2$}.

		Le discriminant est strictement positif, l'équation a deux solutions dans $\R$ qui sont {$\dfrac{4-2}{2}=1$} et {$\dfrac{4+2}{2}=3$}.

		\item $\widehat{\text{DAM}}$ est un angle droit si et seulement si A et M sont deux points distincts vérifiant {$\vectt{AD}\cdot\vectt{AM}=0$}.

		(les points doivent être distincts, sinon A et M ne définissent pas une droite, et donc l'angle $\widehat{\text{DAM}}$ n'est pas défini).

		On connaît déjà les coordonnées de {$\vectt{AD}$}, et celles de {$\vectt{AM}$} sont {$\left(\begin{array}{c}x-3\\ y-1\end{array}\right)$.}

		La condition s'écrit donc \quad {$x-3+y-1=0$}.

		Sachant que \quad {$y=\dfrac3 x$},\quad cette équation devient \quad {$x-3+\dfrac 3 x -1=0$} \quad soit, après multiplication par $x\neq0$ :
		\quad
		{$x^2-4x+3=0$}

		On sait que cette équation a deux solutions : la solution $x=3$ est évidemment sans intérêt : le produit scalaire dans ce cas est nul puisque A et M sont confondus.

		Il reste la solution $x=1$ : le point {M~$(1;3)$} est le seul point de {$\mathcal{H}$} tel que $\widehat{\text{DAM}}$ soit un angle droit.
	\end{enumerate}
